\chapter{常见量子态、非经典性与观测量}
\label{ch:states-observables}

\begin{learningobjectives}
完成本章后，读者应能：
\begin{enumerate}
  \item 写出相干态、热态、压缩高斯态和 Fock 态的 Wigner 函数；
  \item 根据协方差矩阵抽样多模高斯态；
  \item 区分 Wigner 负值、量子性和普通概率可抽样性；
  \item 把等时正规排序观测量转换为 Wigner 样本估计量。
\end{enumerate}
\end{learningobjectives}

\section{选择初态就是选择物理假设}

在 TWA 中，初态不是给平均场加一点任意噪声。初态的粒子数涨落、相位涨落、模间关联和反常关联都由 $\rhohat_0$ 决定。两个具有相同平均粒子数的状态，若一个是相干态、另一个是数态，其长期相位扩散可以完全不同。因此，抽样算法必须从态的矩或协方差出发，而不能只匹配一条平均密度曲线。

\section{相干态与热态}

第2章已经得到相干态
\begin{equation}
  W_{\alpha_0}(\alpha)=\frac{2}{\pi}
  \exp[-2|\alpha-\alpha_0|^2].
\end{equation}
它的位置由 $\alpha_0$ 决定，宽度固定为真空宽度。相干态粒子数服从 Poisson 分布，满足
\begin{equation}
  \qavg{\hat n}=|\alpha_0|^2,
  \qquad
  \Var(\hat n)=|\alpha_0|^2.
\end{equation}
如果实际制备固定总粒子数，直接使用相干态会额外引入全局数涨落；是否可接受取决于目标观测量是否对该涨落敏感。

平均热占据为 $\bar n$ 的单模热态具有
\begin{equation}
  W_{\mathrm{th}}(\alpha)
  =\frac{2}{\pi(2\bar n+1)}
  \exp\!\left[-\frac{2|\alpha|^2}{2\bar n+1}\right].
  \label{eq:thermal-wigner}
\end{equation}
其直接抽样为
\begin{equation}
  \alpha=\sqrt{\frac{\bar n+1/2}{2}}
  (\xi_1+\ii\xi_2),
  \qquad \xi_j\sim\mathcal{N}(0,1).
  \label{eq:thermal-sampling}
\end{equation}
于是 $\ensavg{|\alpha|^2}=\bar n+1/2$，减去排序修正后恢复 $\qavg{\hat n}=\bar n$。

\section{高斯态与协方差矩阵}

连续变量高斯态的辛空间形式、物理性条件与 Williamson 分解可参见综述
\cite{Weedbrook2012}；这里仅保留后续 Wigner 抽样直接需要的结果。

对 $M$ 个模式，把正交分量排列为
\begin{equation}
  \hat{\bm R}=(\hat Q_1,\hat P_1,\ldots,\hat Q_M,\hat P_M)^{\mathsf T}.
\end{equation}
定义均值和对称协方差
\begin{align}
  d_j&=\qavg{\hat R_j},\\
  V_{jk}&=\frac12\qavg{
  \acomm{\hat R_j-d_j}{\hat R_k-d_k}}.
  \label{eq:gaussian-covariance}
\end{align}
采用第2章的无量纲约定 $\comm{\hat Q_j}{\hat P_k}=\ii\delta_{jk}$ 时，物理协方差必须满足
\begin{equation}
  V+\frac{\ii}{2}\Omega\succeq0,
  \qquad
  \Omega=\bigoplus_{j=1}^{M}
  \begin{pmatrix}0&1\\-1&0\end{pmatrix}.
  \label{eq:uncertainty-matrix}
\end{equation}
高斯态的 Wigner 函数就是均值为 $\bm d$、协方差为 $V$ 的普通多元高斯：
\begin{equation}
  W(\bm R)=\frac{1}{(2\pi)^M\sqrt{\det V}}
  \exp\!\left[-\frac12(\bm R-\bm d)^{\mathsf T}
  V^{-1}(\bm R-\bm d)\right].
  \label{eq:multimode-gaussian-wigner}
\end{equation}
若 $V=LL^{\mathsf T}$，则
\begin{equation}
  \bm R=\bm d+L\bm\xi,
  \qquad \bm\xi\sim\mathcal{N}(\bm0,I)
  \label{eq:gaussian-cholesky-sampling}
\end{equation}
给出直接抽样。对合法的有限能量量子协方差，$V$ 是正定矩阵，通常可取 Cholesky 因子。纯高斯态只意味着所有辛本征值等于 $1/2$，并不意味着 $V$ 奇异；真正的数值病态来自强压缩或变量尺度失衡，使普通特征值跨越很多数量级。此时可使用对称特征分解构造平方根，但任何特征值裁剪都必须记录并回测其对目标矩的影响。

\subsection{单模压缩真空}

本书固定约定：$r>0$ 时 $Q$ 正交分量被压缩。此时
\begin{equation}
  V=\frac12
  \begin{pmatrix}
    \ee^{-2r}&0\\
    0&\ee^{2r}
  \end{pmatrix}.
  \label{eq:squeezed-covariance}
\end{equation}
一个正交分量的噪声低于真空，另一个按不确定关系增大。该态具有明确的量子压缩，却仍有处处非负的 Wigner 函数。这再次说明“Wigner 负值”等价于“量子性”的说法过强。

\subsection{模间关联}

两模压缩态的协方差矩阵含非零的模间块。若只按每个模式的边缘方差独立抽样，就会丢掉 $\qavg{\ha_k\ha_{-k}}$ 一类反常关联。淬火后的 Bogoliubov 初态正需要保留这种 $k$ 与 $-k$ 的关联；第15章将给出从准粒子变换构造 $V$ 的方法。

\begin{numericalbox}
高斯抽样至少检查三项：样本均值与 $\bm d$ 的偏差、样本协方差与 $V$ 的矩阵范数误差，以及 $V+\ii\Omega/2$ 的最小特征值。只检查每个对角方差不能发现模间关联丢失。
\end{numericalbox}

\subsection{高斯矩匹配的适用边界}

若目标态本身是高斯态，均值 $\bm d$ 与协方差 $V$ 唯一确定其全部等时矩，\cref{eq:gaussian-cholesky-sampling} 是精确初态抽样。若目标态非高斯，用相同的 $\bm d,V$ 构造高斯分布只保证一、二阶对称矩一致；它不保证粒子数分布、四阶关联、稀有尾部、Wigner 负值或相位干涉一致。

这种差别在非线性演化中不会停留在高阶矩。四次哈密顿量会使低阶矩方程与更高阶矩耦合，所以初始时一、二阶矩相同的两个状态，随后的一、二阶观测量也可能分离。采用高斯矩匹配时必须同时声明：匹配了哪些矩、目标观测量和时间窗口、与精确小系统的比较，以及结论对替代匹配方案是否稳定。不能把“协方差回测通过”写成“目标非高斯态已被精确表示”。

\section{Fock 态与 Wigner 负值}

粒子数态 $\ket{n}$ 的单模 Wigner 函数为
\begin{equation}
  W_n(\alpha)=\frac{2}{\pi}(-1)^n
  L_n(4|\alpha|^2)\ee^{-2|\alpha|^2},
  \label{eq:fock-wigner}
\end{equation}
其中 $L_n$ 是 Laguerre 多项式。除真空外，分布通常含正负交替的环，因此不能把\cref{eq:fock-wigner} 直接当作概率密度生成等权样本。

若 $Z_W=\int\dd^2\alpha\,|W(\alpha)|$ 有限，可定义
\begin{equation}
  P_{|W|}(\alpha)=\frac{|W(\alpha)|}{Z_W},
  \qquad s(\alpha)=\operatorname{sgn}W(\alpha),
\end{equation}
相应的积分估计器为
\begin{equation}
  \int\dd^2\alpha\,W(\alpha)O_W(\alpha)
  =Z_W\mathbb{E}_{|W|}[sO_W].
\end{equation}
这在数学上是无偏的符号重加权，却不是正 Wigner 抽样。平均符号满足
\begin{equation}
  \mathbb{E}_{|W|}[s]=\frac{1}{Z_W};
\end{equation}
当它很小时，所需样本数会急剧增加。更重要的是，这一重加权只处理初态
积分的负权重，并不会补回 TWA 随后舍弃的高阶 Moyal 动力学。

实际工作中有四种选择：在符号问题可控的小系统中使用上述重加权；使用能够处理复权重或扩展相空间的表示；构造只匹配目标低阶矩的近似正分布；或者直接使用精确态作为基准。矩匹配不是精确表示，必须按上一小节说明其边界。

猫态
\begin{equation}
  \ket{\psi_{\mathrm{cat}}}\propto
  \ket{\alpha_0}+\ee^{\ii\phi}\ket{-\alpha_0}
\end{equation}
的 Wigner 函数由两个高斯峰和中间的干涉条纹组成。若只抽样两个峰的经典混合，就会保留“在左或在右”的统计，却丢失相干叠加。Wigner 图中最细的干涉条纹也往往最先暴露半经典截断的局限。

\section{三种常用相空间分布}

不同相空间表示对应不同算符排序：
\begin{center}
\begin{tabular}{lll}
\toprule
表示 & 典型排序 & 主要特征 \\
\midrule
Glauber--Sudarshan $P$ & 正规排序 & 可能比普通函数更奇异 \\
Wigner $W$ & 对称排序 & 实函数，但可为负 \\
Husimi $Q$ & 反正规排序 & 始终非负，但额外平滑 \\
\bottomrule
\end{tabular}
\end{center}
不存在对所有问题都最优的表示。Wigner 表象的优势是漂移方程常接近经典场动力学，且高占据玻色系统的半量子修正清楚；它的限制是非正初态难以直接抽样，非线性演化又会产生高阶导数。
更一般的高斯算符表示可以扩展普通正概率抽样的范围，但其变量、权重和
稳定性条件与本章的 Wigner 高斯抽样不同\cite{Corney2003}。

\section{等时观测量的估计器}

对单模玻色子，常用估计器包括
\begin{align}
  \qavg{\hat n}
  &=\ensavg{|\alpha|^2-\frac12},\\
  \qavg{\hat n(\hat n-1)}
  &=\ensavg{|\alpha|^4-2|\alpha|^2+\frac12},
  \label{eq:wigner-factorial-moments}
\end{align}
从而
\begin{equation}
  g^{(2)}=
  \frac{\qavg{\had\had\ha\ha}}
  {\qavg{\had\ha}^2}.
\end{equation}
对不同模式 $j\neq k$，算符彼此对易，但每个同模的产生湮灭配对仍需排序修正。连续场中的 $1/2$ 会变成含网格尺度的 $1/(2\Delta x)$，第10章将详细推导。

以上规则直接适用于单一时刻的对称排序量。多时关联如 $\qavg{\had(t_1)\ha(t_2)}$ 不能总是把两个轨迹值直接相乘；时间排序、对易子和响应修正可能参与。若目标是谱函数或响应函数，必须从所需的时间排序重新推导估计器。

\begin{warningbox}
状态抽样正确，不代表观测量估计器自动正确。初态分布、动力学方程和算符排序是三项独立输入；任何一项使用了另一种相空间约定，最终结果都可能相差一个随模式数累积的真空修正。
\end{warningbox}

\section*{本章小结}

相干态和热态可以逐模直接抽样，一般高斯态应由完整协方差矩阵抽样。压缩态说明正 Wigner 函数仍能代表量子态，Fock 态和猫态则展示普通概率抽样的限制。等时观测量必须从 Weyl 符号构造；正规排序、多时排序和连续场极限需要分别处理。

\begin{exercises}
  \item 用\cref{eq:thermal-sampling} 计算 $\ensavg{|\alpha|^2}$ 和 $\ensavg{|\alpha|^4}$，验证热态的 $g^{(2)}=2$。
  \item 验证压缩真空协方差\cref{eq:squeezed-covariance} 饱和不确定关系。
  \item 对给定的 $4\times4$ 两模协方差矩阵，写出检查其物理性的步骤和抽样算法。
  \item 计算 $n=1$ Fock 态 Wigner 函数在 $\alpha=0$ 处的值，并解释为什么不能直接作正概率抽样。
  \item 比较相干态与固定粒子数态作为双阱初态时，总粒子数方差的差别，并讨论它可能怎样影响相位扩散。
\end{exercises}
